% ================================================================
%  DOCUMENTAÇÃO DE INTERFACE — SISTEMA DE AGENDAMENTO
%  TCC — UNDB — Guilherme Braga
%  Inserir nos capítulos indicados pelo marcador [CAPÍTULO X]
% ================================================================

% ----------------------------------------------------------------
%  [CAPÍTULO 3 — TECNOLOGIAS UTILIZADAS]
%  Subseção: Tecnologias de Interface e Frontend
% ----------------------------------------------------------------
\subsection{Tecnologias de Interface e \textit{Frontend}}

A camada de apresentação do sistema foi construída inteiramente sobre
tecnologias consolidadas no ecossistema Java, garantindo compatibilidade
nativa com o \textit{framework} Spring Boot, manutenibilidade e facilidade de
evolução futura.

\subsubsection{Thymeleaf}
O Thymeleaf~\cite{thymeleaf} é um motor de \textit{templates} para a plataforma
Java integrado ao Spring MVC. Toda a geração de HTML é realizada no servidor
(\textit{server-side rendering}), eliminando a necessidade de APIs REST separadas
para as telas administrativas e o portal do cliente. Os \textit{templates} herdam
de layouts-base por meio dos mecanismos de fragmento
(\texttt{th:replace} / \texttt{th:fragment}), o que garante consistência visual
e evita duplicação de código.

\subsubsection{Bootstrap 5}
O Bootstrap 5.3.3~\cite{bootstrap5} fornece o sistema de \textit{grid} responsivo
(\textit{Flexbox}-based), os componentes visuais (Navbar, Modal, Card, Badge,
Table) e o conjunto de utilitários CSS utilizados em todas as telas.
A versão 5 eliminou a dependência do jQuery, tornando-o mais leve e adequado ao
JavaScript puro utilizado nos \textit{scripts} de validação e interação.

\subsubsection{Bootstrap Icons}
A biblioteca Bootstrap Icons 1.11.3~\cite{bootstrapicons} é empregada como
sistema de ícones SVG (\textit{font-based}), presente em botões, títulos de página,
itens de menu e \textit{badges} de status. O uso consistente de ícones
contextualiza as ações para o usuário, reduzindo a carga cognitiva.

\subsubsection{Chart.js}
O Chart.js 4.4.3~\cite{chartjs} é incorporado via CDN para renderizar o gráfico
de \textit{doughnut} de distribuição de agendamentos por status no painel
administrativo (\textit{dashboard}). Os dados são injetados pelo Thymeleaf por meio
de atributos \texttt{data-*} em um elemento auxiliar oculto, evitando qualquer
risco de XSS por interpolação direta em JavaScript.


% ----------------------------------------------------------------
%  [CAPÍTULO 4 — DESENVOLVIMENTO DO SISTEMA]
%  Seção: Interface do Sistema
% ----------------------------------------------------------------
\section{Interface do Sistema}

A interface do sistema é dividida em dois ambientes distintos: o
\textbf{Painel Administrativo}, acessível pela equipe interna por meio de
autenticação via Spring Security, e o \textbf{Portal do Cliente}, disponível
publicamente para que clientes realizem agendamentos sem necessidade de cadastro
de usuário administrativo.

% ......................
\subsection{Tela de Login Administrativo}

A tela de login (\texttt{/web/login}) é a porta de entrada do painel
administrativo (Figura~\ref{fig:tela-login}). Seu design prioriza clareza e
segurança visual:

\begin{itemize}
  \item \textbf{Fundo escuro com gradiente} — transmite seriedade e delimita o
    escopo de acesso restrito.
  \item \textbf{Card centralizado} — agrupa os campos de autenticação em um
    contêiner com sombra profunda (\texttt{box-shadow: 0 24px 60px rgba(0,0,0,.4)}),
    destacando o formulário sobre o fundo.
  \item \textbf{Ícones contextuais nos campos} — o ícone de cadeado no campo de
    senha e o ícone de usuário no campo de credencial reduzem a necessidade de
    leitura do rótulo.
  \item \textbf{Alternância de visibilidade da senha} — um botão de olho
    (\textit{password toggle}) ao lado direito do campo de senha permite ao usuário
    visualizar o texto digitado, reduzindo erros de autenticação.
\end{itemize}

\begin{figure}[H]
  \centering
  \includegraphics[width=0.55\textwidth]{figuras/tela-login.png}
  \caption{Tela de login do painel administrativo.}
  \label{fig:tela-login}
\end{figure}

% ......................
\subsection{Painel Administrativo — Navbar e Layout Base}

O layout administrativo (\texttt{layout/base.html}) provê uma barra de navegação
superior fixa (\textit{sticky-top}) com as seguintes características
(Figura~\ref{fig:admin-navbar}):

\begin{itemize}
  \item \textbf{Marca/logo} com ícone \textit{Bootstrap Icons} e nome do sistema.
  \item \textbf{Ícones nos itens de menu} — cada seção (Dashboard, Clientes,
    Profissionais, Serviços, Agendamentos) possui ícone semântico.
  \item \textbf{Responsividade} — colapsa em \textit{hamburger menu} em telas
    menores que 992 px.
  \item \textbf{Atalhos laterais} — botões \textit{Portal} e \textit{Sair} no lado
    direito facilitam a navegação entre os dois ambientes.
  \item \textbf{Destaque do item ativo} — via JavaScript, o link da página
    corrente recebe a classe \texttt{active} automaticamente com base no
    \texttt{window.location.pathname}.
  \item \textbf{Alertas dismissíveis} — mensagens de sucesso e erro exibidas com
    ícone e botão de fechamento automático.
\end{itemize}

\begin{figure}[H]
  \centering
  \includegraphics[width=0.95\textwidth]{figuras/admin-navbar.png}
  \caption{Barra de navegação administrativa com ícones e responsividade.}
  \label{fig:admin-navbar}
\end{figure}

% ......................
\subsection{\textit{Dashboard} — Painel de Controle}

O \textit{dashboard} (\texttt{/dashboard}) é a tela principal do administrador e
concentra as métricas mais relevantes do sistema
(Figura~\ref{fig:dashboard}).

\subsubsection{Cartões de Métricas}
A primeira linha exibe quatro cartões com fundo em gradiente, cada um
representando uma métrica operacional:

\begin{itemize}
  \item \textbf{Total de Agendamentos} — gradiente roxo-lilás.
  \item \textbf{Agendamentos Hoje} — gradiente laranja-dourado.
  \item \textbf{Próximos 7 Dias} — gradiente verde-esmeralda.
  \item \textbf{Clientes com Agendamentos Ativos} — gradiente azul-royal.
\end{itemize}

Cada cartão apresenta o valor numérico em tipografia \textit{bold} de grande
porte, um rótulo descritivo e um ícone temático. O efeito \textit{hover} eleva
o cartão (transform translateY) comunicando interatividade.

\subsubsection{Gráfico de Distribuição por Status}
Um gráfico de \textit{doughnut} (Chart.js) na segunda linha apresenta a
distribuição percentual dos agendamentos por status
(Figura~\ref{fig:dashboard-chart}). O esquema de cores é semanticamente
consistente com os \textit{badges} utilizados nas tabelas:

\begin{table}[H]
\centering
\begin{tabular}{lll}
\hline
\textbf{Status} & \textbf{Cor} & \textbf{Código Hexadecimal} \\ \hline
AGENDADO    & Azul-violeta & \texttt{\#667eea} \\
CONFIRMADO  & Verde        & \texttt{\#28a745} \\
CONCLUÍDO   & Cinza        & \texttt{\#6c757d} \\
CANCELADO   & Vermelho     & \texttt{\#dc3545} \\ \hline
\end{tabular}
\caption{Paleta de cores dos status de agendamento.}
\label{tab:cores-status}
\end{table}

\noindent Os dados são fornecidos pelo \texttt{DashboardWebController} que calcula,
via \texttt{Collectors.groupingBy}, a contagem de cada status e adiciona o
\textit{map} \texttt{porStatus} ao modelo Thymeleaf. O gráfico é alimentado por
atributos \texttt{data-*} em um elemento auxiliar, garantindo separação entre
lógica e apresentação.

\subsubsection{Próximos Agendamentos}
Ao lado do gráfico, uma tabela compacta lista os três próximos agendamentos com
status AGENDADO ou CONFIRMADO, ordenados cronologicamente. Isso permite ao
administrador visualizar rapidamente a agenda sem navegar até a listagem
completa.

\subsubsection{Acesso Rápido aos Cadastros}
A terceira linha exibe três \textit{cards} de acesso rápido — Clientes,
Profissionais e Serviços — com contadores e botões diretos para listagem e
criação de novos registros.

\begin{figure}[H]
  \centering
  \includegraphics[width=0.95\textwidth]{figuras/dashboard.png}
  \caption{\textit{Dashboard} com métricas, gráfico e acesso rápido.}
  \label{fig:dashboard}
\end{figure}

\begin{figure}[H]
  \centering
  \includegraphics[width=0.55\textwidth]{figuras/dashboard-chart.png}
  \caption{Gráfico \textit{doughnut} de distribuição de agendamentos por status.}
  \label{fig:dashboard-chart}
\end{figure}

% ......................
\subsection{Listagens Administrativas}

Todas as telas de listagem (Clientes, Profissionais, Serviços, Agendamentos)
seguem o mesmo padrão visual (Figuras~\ref{fig:clientes-lista},
\ref{fig:agendamentos-lista}):

\begin{itemize}
  \item \textbf{Cabeçalho} com título, contador de registros e botão de criação
    alinhado à direita.
  \item \textbf{Barra de busca/filtro} em cartão separado, com busca em tempo real
    por JavaScript (clientes, profissionais, serviços) ou filtros de formulário
    GET (agendamentos — status e intervalo de datas).
  \item \textbf{Tabela Bootstrap} com classes \texttt{table-hover} e
    \texttt{align-middle}, cabeçalhos em caixa-alta com cor atenuada e tamanho
    reduzido.
  \item \textbf{Botões de ação} ícone-apenas (\textit{outline}, tamanho sm) para
    \textit{Ver}, \textit{Editar}, \textit{Alterar Status} e \textit{Deletar}.
  \item \textbf{Badges de status} com cores semânticas (amarelo, verde, cinza,
    vermelho) para indicar o estado de cada agendamento.
  \item \textbf{Estado vazio} com ícone e \textit{call to action} para criar o
    primeiro registro.
  \item \textbf{Modal de confirmação} centralizado para operações destrutivas,
    exibindo os dados do registro e alerta irreversível antes da exclusão.
\end{itemize}

\begin{figure}[H]
  \centering
  \includegraphics[width=0.95\textwidth]{figuras/clientes-lista.png}
  \caption{Listagem de clientes com busca e ações por ícone.}
  \label{fig:clientes-lista}
\end{figure}

\begin{figure}[H]
  \centering
  \includegraphics[width=0.95\textwidth]{figuras/agendamentos-lista.png}
  \caption{Listagem de agendamentos com filtros e badges de status.}
  \label{fig:agendamentos-lista}
\end{figure}

% ......................
\subsection{Formulários Administrativos}

Os formulários de cadastro e edição (Figuras~\ref{fig:form-cliente},
\ref{fig:form-agendamento}) são apresentados em um contêiner \texttt{form-card}
de largura máxima de 640 px, centralizado no conteúdo principal. Os campos
utilizam \textit{labels} explícitos, marcação de obrigatoriedade com asterisco
vermelho, \textit{placeholders} com exemplos, validação nativa HTML5 e
\textit{feedback} de erro integrado ao ciclo de \textit{binding} do Spring MVC
(classe \texttt{is-invalid} + \texttt{invalid-feedback} do Bootstrap).

\begin{figure}[H]
  \centering
  \includegraphics[width=0.65\textwidth]{figuras/form-cliente.png}
  \caption{Formulário de cadastro de cliente com validação de campos.}
  \label{fig:form-cliente}
\end{figure}

\begin{figure}[H]
  \centering
  \includegraphics[width=0.65\textwidth]{figuras/form-agendamento.png}
  \caption{Formulário de novo agendamento com info dinâmica do serviço.}
  \label{fig:form-agendamento}
\end{figure}

% ......................
\subsection{Portal do Cliente}

O portal público (\texttt{/portal}) permite ao cliente navegar pelos serviços
disponíveis, realizar agendamentos e acompanhar seu histórico, sem acesso ao
painel administrativo.

\subsubsection{\textit{Home} do Portal}
A página inicial (\texttt{/portal}, Figura~\ref{fig:portal-home}) exibe:

\begin{itemize}
  \item \textbf{\textit{Hero card}} com gradiente violeta-roxo, título de
    boas-vindas, descrição e dois botões de ação principal (Ver Serviços e
    Agendar Agora). Um ícone decorativo grande preenchido com baixa opacidade
    ocupa o lado direito em telas médias/grandes, criando profundidade visual.
  \item \textbf{Estatísticas em cartões}: Serviços Disponíveis, Profissionais e
    Meus Agendamentos, cada um com ícone em plano de cor temático, valor
    numérico destacado e rótulo em caixa-alta.
\end{itemize}

\begin{figure}[H]
  \centering
  \includegraphics[width=0.95\textwidth]{figuras/portal-home.png}
  \caption{\textit{Home} do portal do cliente com \textit{hero} e estatísticas.}
  \label{fig:portal-home}
\end{figure}

\subsubsection{Navbar do Portal}
A barra de navegação do portal utiliza gradiente primário
(\texttt{linear-gradient(135deg, \#667eea 0\%, \#764ba2 100\%)}), logomarca com
ícone de calendário, links com destaque ativo e, quando o cliente está
autenticado, exibição do nome e botão de Sair.

% ......................
\subsection{\textit{Wizard} de Agendamento}

O fluxo de agendamento é conduzido por um assistente em quatro etapas
(\textit{wizard}), acessível em \texttt{/portal/step1} a \texttt{/portal/step4}.
O indicador de progresso no topo de cada página comunica visualmente em qual
etapa o cliente se encontra (Figura~\ref{fig:wizard-steps}).

\begin{description}
  \item[Etapa 1 — Serviço] Seleção do serviço por meio de \textit{cards}
    clicáveis em \textit{grid} responsivo (col-md-6, col-lg-4). Cada
    \textit{card} exibe nome, descrição, preço e duração. O \textit{card}
    selecionado recebe borda de destaque e fundo levemente colorido.

  \item[Etapa 2 — Data e Hora] Seleção de profissional, data (campo
    \texttt{date} com data mínima = hoje) e horário (botões de
    \textit{time-slot} gerados dinamicamente em intervalos de 30 minutos das
    8h às 18h). A etapa anterior aparece marcada com ícone de verificação.

  \item[Etapa 3 — Confirmação] Resumo completo do agendamento (serviço,
    profissional, data, horário, valor) com alerta informativo sobre o status
    inicial e botão de confirmação.

  \item[Etapa 4 — Concluído] Tela de sucesso com ícone verde, número de
    protocolo do agendamento e resumo final. Oferece navegação para \textit{Meus
    Agendamentos} ou novo agendamento.
\end{description}

\begin{figure}[H]
  \centering
  \includegraphics[width=0.75\textwidth]{figuras/wizard-steps.png}
  \caption{Indicador de progresso e \textit{cards} de serviço da Etapa 1.}
  \label{fig:wizard-steps}
\end{figure}

\begin{figure}[H]
  \centering
  \includegraphics[width=0.65\textwidth]{figuras/wizard-confirmacao.png}
  \caption{Etapa 3 do \textit{wizard}: revisão e confirmação do agendamento.}
  \label{fig:wizard-confirmacao}
\end{figure}

\begin{figure}[H]
  \centering
  \includegraphics[width=0.65\textwidth]{figuras/wizard-sucesso.png}
  \caption{Etapa 4: tela de sucesso com protocolo de confirmação.}
  \label{fig:wizard-sucesso}
\end{figure}


% ----------------------------------------------------------------
%  [CAPÍTULO 4 — DESENVOLVIMENTO DO SISTEMA]
%  Seção: Melhorias de Usabilidade Implementadas
% ----------------------------------------------------------------
\section{Melhorias de Usabilidade Implementadas}

Durante o desenvolvimento do sistema, foram identificadas e corrigidas diversas
deficiências de usabilidade nas interfaces, alinhando o produto às heurísticas
de Nielsen~\cite{nielsen1994} e às diretrizes WCAG 2.1~\cite{wcag21}.

\subsection{Consistência Visual (\textit{Design System})}

Toda a aplicação compartilha um conjunto de variáveis CSS definidas em
\texttt{style.css} (\texttt{:root}), abrangendo paleta de cores, escala de
espaçamento, raios de borda e sombras. Isso garante que qualquer alteração de
identidade visual seja realizada em um único local e propagada
automaticamente a todos os componentes.

\subsection{Navegação com Ícones}

Os menus de navegação — tanto do painel administrativo quanto do portal —
incorporam ícones do Bootstrap Icons ao lado dos rótulos textuais. Estudos de
usabilidade demonstram que ícones associados a texto aceleram o reconhecimento
de destinos de navegação em comparação ao texto puro~\cite{pierce2004}.

\subsection{Alternância de Visibilidade de Senha}

O campo de senha na tela de login passou a exibir um botão de olho
(\textit{toggle}), permitindo ao usuário alternar entre texto oculto e visível.
Essa funcionalidade reduz erros de digitação e é recomendada pelas diretrizes de
autenticação do NIST SP 800-63B~\cite{nist80063b}.

\subsection{Tabelas Aprimoradas}

As tabelas de Profissionais e Serviços foram inteiramente reescritas para utilizar
as classes Bootstrap (\texttt{table-hover}, \texttt{align-middle},
\texttt{table-light} no cabeçalho). As ações passaram de textos simples para
botões ícone-apenas com \textit{tooltips}, reduzindo a poluição visual sem perder
clareza de função. Modais de confirmação substituem os \texttt{window.confirm()}
nativos do navegador, mantendo consistência de UX.

\subsection{Correção de Defeitos HTML}

A listagem de agendamentos apresentava estrutura HTML malformada (elementos
\texttt{<table>} sem classes Bootstrap, \texttt{</div>} órfãos e referências a
modais inexistentes). A tela foi reescrita com HTML válido, preservando toda a
lógica de transição de status e exclusão.

\subsection{Portal do Cliente — \textit{Hero Card}}

A página inicial do portal foi aprimorada com um cartão \textit{hero} em
gradiente, substituindo o layout simples anterior. O contraste entre o fundo
gradiente e os botões em estilo claro (\texttt{btn-light} e
\texttt{btn-outline-light}) segue as recomendações de contraste da WCAG 2.1
(taxa mínima de 4.5:1 para texto normal).

\subsection{Gráfico de Pizza — Monitoramento de Métricas}

A adição do gráfico de \textit{doughnut} no \textit{dashboard} proporciona uma
visão instantânea da distribuição do fluxo de atendimentos, permitindo ao
gestor identificar gargalos (excesso de agendamentos não confirmados, por
exemplo) sem necessidade de consultar a listagem completa.

\subsection{Fluxo de Wizard — Remoção de Scripts Duplicados}

A Etapa 1 do \textit{wizard} de agendamento continha dois blocos
\texttt{<script>} que registravam os mesmos \textit{event listeners}, causando
execução duplicada das funções de seleção de serviço. O bloco redundante foi
removido, mantendo apenas a implementação otimizada que utiliza referência ao
cartão ativo (O(1)) em vez de iterar todos os cartões a cada clique.


% ----------------------------------------------------------------
%  REFERÊNCIAS — adicionar ao arquivo .bib
% ----------------------------------------------------------------
%
% @misc{thymeleaf,
%   author  = {{Thymeleaf Team}},
%   title   = {Thymeleaf 3.x Documentation},
%   year    = {2024},
%   url     = {https://www.thymeleaf.org/documentation.html}
% }
% @misc{bootstrap5,
%   author  = {{The Bootstrap Authors}},
%   title   = {Bootstrap 5.3 Documentation},
%   year    = {2024},
%   url     = {https://getbootstrap.com/docs/5.3}
% }
% @misc{bootstrapicons,
%   author  = {{The Bootstrap Authors}},
%   title   = {Bootstrap Icons 1.11},
%   year    = {2024},
%   url     = {https://icons.getbootstrap.com}
% }
% @misc{chartjs,
%   author  = {{Chart.js Contributors}},
%   title   = {Chart.js 4.x Documentation},
%   year    = {2024},
%   url     = {https://www.chartjs.org/docs/latest}
% }
% @book{nielsen1994,
%   author    = {Nielsen, Jakob},
%   title     = {Usability Engineering},
%   publisher = {Morgan Kaufmann},
%   year      = {1994}
% }
% @techreport{wcag21,
%   author      = {{W3C Web Content Accessibility Guidelines Working Group}},
%   title       = {{WCAG} 2.1},
%   institution = {W3C},
%   year        = {2018},
%   url         = {https://www.w3.org/TR/WCAG21}
% }
% @techreport{nist80063b,
%   author      = {Grassi, Paul A. and others},
%   title       = {{NIST SP 800-63B}: Digital Identity Guidelines — Authentication},
%   institution = {NIST},
%   year        = {2017}
% }
% @article{pierce2004,
%   author  = {Pierce, Jon L.},
%   title   = {The Role of Icons in User Interface Design},
%   journal = {Journal of Usability Studies},
%   year    = {2004}
% }
